\begin{papercard}{xie2025}
\cf{Stages} \stagetag{L} (open-weight LLM pipeline for clause detection in policies; prompt design from legal text; segmentation; deterministic decoding), \stagetag{C} (a clause catalogue derived from ten privacy laws as codebook), \stagetag{I} (per-clause F1 against mapped human-annotated datasets). USENIX Security 2025 paper (Georgia Tech). Added from the gap search of 23 October 2026 as the most recent large-scale LLM study of policies under GDPR-era law.

\cf{Aim and frame} Prior automated policy analysis relied on rules or supervised classifiers built before the current privacy laws; the authors systematise the practices required by ten laws in force in the EU and the US \pc{xie2025}{2}{We first systematize the privacy practices discussed by 10 notable privacy regulations currently in effect in the European Union and the US, identifying 34 distinct clauses on privacy practices across 4 overarching themes}, build an LLM method to detect them, and apply it to policies collected from more than 100,000 websites \pc{xie2025}{2}{we collect privacy policies from over 100K websites, and apply our LLM method to a subset of sites}.

\cf{Model and configuration} Open-weight Llama-3.1-70B-Instruct, chosen over the 405B model for hardware cost and reproducibility \pc{xie2025}{6}{In this work, we use the Llama-3.1-70B-Instruct [6] as the foundational model for interpreting privacy policies}; run locally in bfloat16 with sampling disabled for deterministic output \pc{xie2025}{8}{We set the max token length to 2048 and disable sampling (do_sample=false) to obtain deterministic results}. Two task types: \emph{coverage} (binary: does a segment contain a clause, plus extraction of the relevant text) and \emph{inspection} (multi-class: which data types, purposes, recipients).

\cf{Prompting} Prompts were written by hand from the wording of the laws, merging the terms different laws use for the same clause. Zero-shot was kept because examples and chain-of-thought did not help \pc{xie2025}{8}{While we evaluated few-shot learning and chain- of-thought [40,65] prompting, common LLM prompting tech- niques, we observed limited improvements or worse perfor- mance on some tasks}.

\cf{Segmentation} Whole documents lowered recall, so policies were converted to Markdown (keeping tables) and split by headings into pieces of at most 1,000 tokens \pc{xie2025}{8}{We also observe in our experiments that inputting the full privacy policy docu- ment often leads to false negatives}; Markdown kept the table context of CCPA-style disclosures \pc{xie2025}{10}{using the Markdown format as LLM inputs helps maintain the table context and the overall structure of the privacy policy}.

\cf{Evaluation data} Labels came from three human-annotated datasets (OPP-115, PPGDPR, C3PA) mapped to the clauses, with the majority label where annotators differed \pc{xie2025}{7}{To determine a final label in OPP-115 and C3PA, we apply the majority rule}. Clauses not covered were labelled by the authors until each had 50 positive and 50 negative segments \pc{xie2025}{7}{we continuously randomly sampled and labeled segments until 50 positive and 50 negative samples were found}. All test sets are balanced between positives and negatives \pc{xie2025}{8}{along with the evaluation dataset used and the total number of test samples (balanced between pos- itive and negative samples}.

\cf{Results} Mean F1 0.94 over the 34 clauses, none below 0.84, against mostly below 0.8 for earlier methods \pc{xie2025}{3}{our LLM approach exhibits an aver- age F1-score of 0.94 (with no task below 0.84)}. Policy discovery itself needed manual checks: 76\,\% accuracy in the ``plausible'' group and 99\,\% in the high-confidence group, which alone entered the measurement \pc{xie2025}{11}{We then manually checked 100 randomly selected domains in the high-confidence group, reaching 99\% accuracy}. Regional differences were large, e.g.\ the CCPA right to non-discrimination appears in 11\,\% of top EU and 61\,\% of top US domains \pc{xie2025}{13}{R10 (right to non-discrimination) is CCPA-specific and is in the policies of only 11\% of top EU do- mains compared to 61\% of top US domains}.

\cf{Limits (card)} Balanced test sets give precision at 50\,\% prevalence; at the real, much lower prevalence of rare clauses precision would be lower, and the paper reports no prevalence-weighted figure. Ground truth for half of the clauses was labelled by the authors without a reported agreement statistic. English website policies only, as the authors state \pc{xie2025}{11}{Our work focuses on systematizing Western privacy laws (Section 3) and analyzing English privacy policies (Section 5), which may not generalize to other regions}. Clause presence is a coarser task than the thesis's harm, stance and protection spans.

\cf{Use for this thesis} (1) A recent, citable configuration for D9: an open-weight model run locally, sampling off, zero-shot prompts written from the codebook's legal sources, Markdown input and segments of about 1,000 tokens. (2) Grounds to test, not assume, that examples help (compare Rodriguez et al.\ \pp{rodriguez2024}, where definitions and two examples helped). (3) The thesis's whole-app validation set is the stronger design; report precision at observed prevalence and say why it differs from balanced benchmarks. (4) The GDPR clause list (rights, contact, retention, recipients) can seed the ``additional details'' and component codes of Axis A.
\end{papercard}
